\section{The original balanced Mellin kernel: full polynomial comparison and meromorphic remainder}
\label{BK:sec:original-balanced-kernel}

The authoritative shared transcript constructs the raw Mellin realization
in A1167, retains its kernel complex in A1361, and replaces it by the
$\mathbb C[t]$-balanced realization in A1427.  This section computes
the polynomial action on the remaining kernel on those original objects.
It uses no replacement arithmetic function and no auxiliary zero packet.

\subsection{Original objects and full-plane Euler inverses}

Retain, exactly as in A1167 and A1427,
\[
 \begin{split}
 V={}&\{\phi\in\mathcal S(\mathbb R):\phi(-x)=\phi(x),\
                 \phi(0)=0,\ \int_{\mathbb R}\phi(x)\,dx=0\},\\
 \mathscr B={}&\{F\in C^\infty(\mathbb R_{>0}):
       \sup_{x>0}x^b|D^kF(x)|<\infty\quad(b\in\mathbb Z,k\ge0)\},
 \qquad D=-x\partial_x,\\
 \Theta\phi(x)={}&2\sum_{n\ge1}\phi(nx),\qquad
 \mathcal MF(s)=\int_0^\infty F(x)x^s\frac{dx}{x},
 \qquad M_+\phi(s)=\int_0^\infty\phi(x)x^s\frac{dx}{x},\\
 g(s)={}&2\xi(s)=s(s-1)\pi^{-s/2}\Gamma(s/2)\zeta(s),\\
 H_\phi(s)={}&\frac{2\pi^{s/2}M_+\phi(s)}
                         {s(s-1)\Gamma(s/2)}.
 \end{split}                                                    \tag{BK.1}
\]
The original Fourier convention is
$\widehat\phi(\nu)=\int_{\mathbb R}\phi(x)e^{-2\pi ix\nu}\,dx$.
The original theta construction proves
\[
 \mathcal M\Theta\phi=gH_\phi,\quad
 \Theta D=D\Theta,\quad
 \widehat{D\phi}=(1-D)\widehat\phi,\quad
 H_{D\phi}(s)=sH_\phi(s),\quad
 H_{\widehat\phi}(s)=H_\phi(1-s).                 \tag{BK.2}
\]
The last equality follows directly from Poisson's identity
$\Theta\widehat\phi(x)=x^{-1}\Theta\phi(1/x)$, Mellin substitution,
and $g(1-s)=g(s)$; cancellation holds on the nonempty open set
$g\ne0$ and then throughout the plane by holomorphy.
Here $H_\phi$ is entire: Taylor subtraction of the even smooth
function at zero gives only simple possible poles of $M_+\phi$
at the negative even integers; the corresponding zeros of
$1/\Gamma(s/2)$ cancel them.  At $s=0$, $M_+\phi$ is holomorphic
because $\phi(0)=0$, and $s\Gamma(s/2)$ tends to $2$.
At $s=1$, $M_+\phi(1)=0$ is exactly the retained integral condition.
In particular
\[
 H_\phi(1)=2(M_+\phi)'(1).                       \tag{BK.3}
\]
The topology and all these maps are the original Schwartz ones.

For every $\rho\in\mathbb C$ one has the exact range formulas
\[
 (D-\rho)\mathscr B=\ker(\mathcal M(\cdot)(\rho)),\qquad
 (D-\rho)V=\ker(H_{(\cdot)}(\rho)).             \tag{BK.4}
\]
Both operators are injective.  We now prove the full-plane source
formula, including the endpoint $\rho=1$ that cannot be handled by
dividing by $1-\rho$.

On the target zero-Mellin domain the inverse is the original integral
\[
 S_\rho F(x)=x^{-\rho}\int_x^\infty F(y)y^{\rho-1}\,dy.
                                                               \tag{BK.5}
\]
Differentiating gives $(D-\rho)S_\rho F=F$.  Since the total Mellin
moment is zero, the same expression is
$-x^{-\rho}\int_0^xF(y)y^{\rho-1}dy$ near zero.  For
$\sigma=\Re\rho$, an upper bound $|F(y)|\le C_Ny^{-N}$ yields
$|S_\rho F(x)|\le C_Nx^{-N}/(N-\sigma)$ when $N>\sigma$.
A lower bound $|F(y)|\le C'_Ny^N$ yields
$|S_\rho F(x)|\le C'_Nx^N/(N+\sigma)$ when $N+\sigma>0$.
The differential identity controls each further Euler derivative.
Thus every original seminorm is finite and the inverse is continuous
on the indicated kernel.  The only homogeneous solution is
$cx^{-\rho}$, which belongs to $\mathscr B$ only when $c=0$.
Integration by parts proves
$\mathcal M((D-\rho)F)=(s-\rho)\mathcal MF$, so the range is
exactly the first kernel in (BK.4).

For the source, first let $\Re\rho>0$ and $H_\phi(\rho)=0$.
For $\rho\ne1$, the multiplier of $M_+\phi(\rho)$ in (BK.1)
is finite and nonzero, so this Mellin moment vanishes.  For $\rho=1$
the same moment already vanishes by the definition of $V$.
The integral (BK.5) is consequently also
\[
 S_\rho\phi(x)=-\int_0^1u^{\rho-1}\phi(xu)\,du. \tag{BK.6}
\]
For every ordinary derivative of order $k$, differentiating this
integral inserts $u^k\phi^{(k)}(xu)$.  Uniform domination on bounded
$x$ intervals follows from $\Re\rho+k>0$.  Therefore the result
extends smoothly and evenly across zero with value zero there.
The upper-tail expression gives Schwartz decay, with every derivative,
at infinity.  If $\rho\ne1$, integration of
$(D-\rho)S_\rho\phi=\phi$ over $\mathbb R$ gives
\[
 (1-\rho)\int_{\mathbb R}S_\rho\phi=\int_{\mathbb R}\phi=0.
\]
If $\rho=1$, integration by parts in a right half-plane gives
\[
 M_+(S_1\phi)(s)=\frac{M_+\phi(s)}{s-1}.
\]
Taking its removable value at $1$ and using (BK.3) gives
$\int_{\mathbb R}S_1\phi=2(M_+\phi)'(1)=H_\phi(1)=0$.
Thus the original source integral condition is retained at every
$\rho$ with positive real part.

If $\Re\rho\le0$, put $\lambda=1-\rho$, so $\Re\lambda\ge1$.
Fourier preserves $V$, and (BK.2) gives
$H_{\widehat\phi}(\lambda)=H_\phi(\rho)=0$.  The exact inverse is
\[
 S_\rho^V\phi=-\widehat{S_\lambda\widehat\phi}. \tag{BK.7}
\]
On even functions Fourier squares to the identity.  The intertwiner
in (BK.2) therefore gives
$(D-\rho)(-\widehat{S_\lambda\widehat\phi})=\phi$,
with every source condition preserved.  For positive real part the
inverse is (BK.6).  The homogeneous solution cannot be Schwartz at
infinity unless zero, so the inverse is unique.  Identity (BK.2)
proves the reverse containment in the second formula of (BK.4).
This completes the full-plane calculation without excluding either
endpoint or negative even Taylor coordinate.

\subsection{Every polynomial kernel and cokernel is computed on the actual theta sequence}

Set $A=\mathbb C[t]$, with $t$ acting as the original $D$, and
$Q=\mathscr B/\Theta V$.  For a nonconstant monic polynomial
\[
 h(t)=\prod_\rho(t-\rho)^{n_\rho},\qquad
 E_h=A/(h),
\]
let $j_h$ retain all Taylor jets of orders less than $n_\rho$.
Iterating (BK.4) gives
\[
 h(D)V=\ker(j_hH),\qquad
 h(D)\mathscr B=\ker(j_h\mathcal M).             \tag{BK.8}
\]
Indeed after dividing by $D-\rho$, the Mellin multiplier or $H$
is divided by $s-\rho$ by (BK.2); the remaining zero orders are
exactly the orders needed at the next division.  Injectivity proves
that the iterated inverse is independent of the chosen ordering.

The source jet map is onto because the original vector
\[
 \phi_*=(4\pi^2x^4-6\pi x^2)e^{-\pi x^2}
 \quad\text{satisfies}\quad H_{P(D)\phi_*}=P.    \tag{BK.9}
\]
The target jet map is onto because its functionals on compactly
supported smooth logarithmic functions are the independent functions
$y^j e^{\rho y}/j!$.  For precision, a dependence vanishing on
an interval extends to an entire exponential polynomial.  Applying
$\prod_{\sigma\ne\rho}(\partial_y-\sigma)^{n_\sigma}$ removes
the other summands; its restriction to $e^{\rho y}$ times
polynomials of degree less than $n_\rho$ is invertible, since it
is a product of nonzero constants plus nilpotent derivatives.
Then $(\partial_y-\rho)$ repeatedly separates the polynomial
coefficients.  All coefficients vanish.  If the finite vector-valued
jet map were not onto, a nonzero linear functional on its missing
quotient would give such a dependence, a contradiction.

Consequently the quotient square and exact theta sequence are
\[
 \begin{array}{ccc}
 V/h(D)V&\xrightarrow{\Theta}&\mathscr B/h(D)\mathscr B\\
 j_hH\downarrow\wr&&\wr\downarrow j_h\mathcal M\\
 E_h&\xrightarrow{M_{j_hg}}&E_h,
 \end{array}                                                    \tag{BK.10}
\]
\[
 0\longrightarrow Q[h(D)]\xrightarrow{\delta_h}E_h
 \xrightarrow{M_{j_hg}}E_h\longrightarrow Q/h(D)Q\longrightarrow0.
                                                               \tag{BK.11}
\]
These identities now hold for every polynomial $h$, without a
strip restriction.  The map $\delta_h$ sends $[F]$ to $j_hH_\phi$
where $h(D)F=\Theta\phi$.  The source $\phi$ is unique because
$\Theta$ is injective.  Changing $F$ by a theta relation changes
$\phi$ by an $h(D)$ multiple, so the map is well-defined.
If its value is zero, (BK.8) writes $\phi=h(D)\psi$ and
injectivity of $h(D)$ on $\mathscr B$ gives $F=\Theta\psi$.
A vector in the kernel of $M_{j_hg}$ is represented by a polynomial
$P$; the zero-jet condition permits applying the actual inverse of
$h(D)$ to $\Theta(P(D)\phi_*)$, producing its preimage.
Quotienting the right target in (BK.10) proves the last cokernel.
This proves the entire sequence with its original maps.

\subsection{The balanced kernel has invertible polynomial action}

Fix any analytic stalk $O=\mathcal O_{\mathbb C,s_0}$, with
$t$ acting on $O$ by multiplication by $s$.  Write
\[
 M=O\otimes_AQ,\qquad N=O/(g),\qquad
 b:M\to N,\quad b(a\otimes[F])=[a\mathcal MF],\qquad K=\ker b.
                                                               \tag{BK.12}
\]
The map is balanced by (BK.2).  The ring $O$ is torsion-free over
the PID $A$: a nonzero polynomial times a holomorphic germ is zero
only when that germ is zero.  It is therefore flat.  One direct
justification of this familiar PID fact is to write a torsion-free
module as the directed union of its finitely generated submodules;
each such submodule is free by the PID structure theorem.  Tensoring
with a free module preserves injections, and filtered colimits
preserve exactness, proving flatness.

The map $b$ is onto.  In fact the original test function
\[
 F_0(x)=\exp[-(\log x)^2],\qquad
 H_0(s)=\mathcal MF_0(s)=\sqrt\pi\exp(s^2/4)     \tag{BK.13}
\]
belongs to $\mathscr B$ and $H_0$ is an invertible germ at every
point.  Euler derivatives of $F_0$ are polynomials in $\log x$
times $F_0$, and $x^bF_0$ is bounded for every integer $b$.
Its Mellin integral becomes
$\int_{\mathbb R}e^{-y^2+sy}dy$; real completion of the square
gives the formula for real $s$, and entire continuation gives it
for every complex $s$.  Thus
$b(aH_0^{-1}\otimes[F_0])=[a]$, proving surjectivity.

For every nonzero polynomial $h$, the two induced maps are
isomorphisms:
\[
 b:M[h]\xrightarrow{\sim}N[h],\qquad
 \overline b:M/hM\xrightarrow{\sim}N/hN.        \tag{BK.14}
\]
For a nonzero constant both statements are immediate.  For a
nonconstant polynomial flatness gives
$M[h]=O\otimes_AQ[h(D)]$ and
$M/hM=O\otimes_A(Q/h(D)Q)$.
Also $O\otimes_AE_h=O/(h)$ by the explicit map
$a\otimes[P(t)]\mapsto[aP(s)]$, with inverse $[a]\mapsto a\otimes1$.
The cokernel assertion in (BK.11) consequently gives
\[
 M/hM\cong O/(h,g)=N/hN,
\]
and the isomorphism is precisely the Mellin map $\overline b$.

For the kernel map retain the actual local units.  Put $z=s-s_0$
and write
\[
 h(s)=a(z)z^n,\qquad g(s)=u(z)z^m,
 \qquad a(0)u(0)\ne0,\quad n,m\ge0.             \tag{BK.15}
\]
Under (BK.11), a local vector $v$ in
$\ker(g:O/(h)\to O/(h))$ corresponds to a torsion class with
$h\mathcal MF=gH_\phi$, and its image under $b$ is exactly
\[
 v\longmapsto\left[\frac gh\,v\right]\in N[h]. \tag{BK.16}
\]
This is well-defined: changing $v$ by $h$ times a germ changes
the displayed product by $g$ times a germ.  The kernel condition
makes the product holomorphic.  Explicitly its domain is
$z^{\max(n-m,0)}O/z^nO$ and its target is
$z^{\max(m-n,0)}O/z^mO$.  Formula (BK.16) is multiplication by
the actual unit $u/a$, together with the power $z^{m-n}$.
It sends precisely the indicated source powers to the indicated
target powers and has inverse multiplication by $a/u$ and the
power $z^{n-m}$.  It is an isomorphism, including $m=0$ or $n=0$
when both spaces are zero.  This proves the first map in (BK.14)
without dropping a local factor of $g$ or a repeated-zero order.

Apply the kernel-cokernel exact sequence for multiplication by $h$
to $0\to K\to M\xrightarrow bN\to0$.  It is
\[
 0\to K[h]\to M[h]\to N[h]\to K/hK\to M/hM\to N/hN\to0.
                                                               \tag{BK.17}
\]
Both maps in (BK.14) are isomorphisms, so $K[h]=0$ and $K/hK=0$.
We have proved
\[
 \boxed{\quad h:K\xrightarrow{\sim}K
             \quad\text{for every nonzero }h\in\mathbb C[t].\quad}
                                                               \tag{BK.18}
\]
This is a computed property of the original balanced kernel.

It has a stronger local form.  Every nonzero holomorphic germ
is a unit times $z^r$.  Units already act bijectively, and (BK.18)
applied to $(t-s_0)^r$ makes $z^r$ bijective.  Hence $K$ has a
unique module structure over the meromorphic germ field
$\operatorname{Frac}O$, obtained by taking these actual inverses:
\[
 \frac ab\cdot k=a\,(b|_K)^{-1}k.             \tag{BK.19}
\]
If $a/b=c/d$, the equality $ad=bc$ and invertibility of $bd$
show that the two prescribed actions agree.  The field laws follow
by multiplying through their nonzero denominators.

\subsection{Canonical finite spectral part and an explicit nonzero kernel class}

At a point where $m=\operatorname{ord}_{s_0}g>0$, (BK.14) for
$h=(t-s_0)^m$ gives a unique inverse
$j:N\xrightarrow{\sim}M[z^m]$ of $b$ on that kernel.  Therefore
\[
 M=M[z^m]\oplus K,\qquad
 x=j(bx)+(x-j(bx)),\qquad M[z^m]\cong O/(g).   \tag{BK.20}
\]
The first summand is all the $O$-torsion of $M$.  Indeed every
torsion element is killed by a power of $z$.  For larger powers
than $m$, (BK.14) still gives the same unique preimage of each
element of $N$, because $K$ has no $z$-torsion.  When $m=0$,
$N=0$ and the formula is $M=K$.  All the full nilpotent jets of
an actual zero occur in the first summand of (BK.20).

The kernel $K$ is nonzero.  The only classical zero-set input needed
for the following assertion is the already-retained Hardy theorem:
there are distinct $T_j\to+\infty$ with
$g(1/2+iT_j)=0$.  The cumulative reader uses this exact input in
\texttt{theta\_gamma\_reference.tex}, (TG.26), and
\texttt{sum\_connection\_stieltjes.tex}, (FC.21), with Hardy's
original 1914 source.  No assumption of RH is used.

Let $q_0=[F_0]\in Q$.  It is $A$-nontorsion.  If
$P(D)F_0=\Theta\phi$ with a polynomial $P$, then (BK.2) and
(BK.13) give
\[
 P(s)\sqrt\pi e^{s^2/4}=g(s)H_\phi(s).
\]
At every $1/2+iT_j$, the nonzero exponential forces $P(1/2+iT_j)=0$.
A nonzero polynomial cannot have infinitely many distinct roots,
so $P=0$.  Thus $A\to Q$, $P\mapsto P(D)q_0$, is injective.
Flatness now proves that $O\to M$, $a\mapsto a\otimes q_0$,
is injective.  In particular
\[
 0\ne g\otimes q_0\in K.                       \tag{BK.21}
\]

The corresponding cycle is explicit on the original balanced complex.
Put $b_*=\Theta\phi_*$, so $\mathcal Mb_*=g$.  In
$O\otimes_A\mathscr B$, set
\[
 Z_0=g\otimes F_0-H_0\otimes b_*.
 \qquad \beta Z_0=gH_0-H_0g=0.                 \tag{BK.22}
\]
Its degree-one class is exactly the nonzero class in (BK.21),
since its second term is the original theta boundary.  For the
original invariant source label $W_*=\mathbb C[D]\phi_*$,
the map $O\otimes_AW_*\to O$ induced by $H$ is an isomorphism
by (BK.9).  Thus the kernel complex at that label has degree-zero
term zero, and (BK.22) is an explicitly nonzero degree-one kernel
cycle.  Passing to the full label $V$ preserves its nonzero class,
as proved by the $A$-nontorsion argument.  This resolves the status
left open in A1427: balancing removes the earlier displayed
spurious relation, but its full degree-one kernel is not zero.

This statement does not create a spectral obstruction to RH.
For every finite-dimensional $A$-module $E$, let $p_E$ be the
characteristic polynomial of its $t$ action.  Cayley--Hamilton
gives $p_EE=0$.  If $u:E\to K$ is $A$-linear, then
$p_Eu(E)=0$ and injectivity in (BK.18) gives $u=0$.
If $v:K\to E$ is $A$-linear, every $k\in K$ is $p_Ek'$ for
some $k'$, and $v(k)=p_Ev(k')=0$.  Consequently
\[
 \operatorname{Hom}_A(E,K)=0=
 \operatorname{Hom}_A(K,E)\quad(\dim_{\mathbb C}E<\infty).
                                                               \tag{BK.23}
\]
The entire kernel has no finite-dimensional invariant subspace
and no nonzero finite-dimensional equivariant quotient.  This does
not assert the same for arbitrary subquotients of its proper
submodules.  Its calculated contribution to every full finite
spectral jet observation is zero, whereas (BK.20) proves that
the actual finite spectral torsion survives bijectively.

Finally every construction occurs in the original supported quotient.
The structure assertions (BK.18)--(BK.23) concern the full source
label $V$ used in (BK.12).  They are not asserted for an arbitrary
proper $D$-invariant source label.  For such a label $W\subset V$,
the exact comparison that must be retained is
\[
 0\longrightarrow V/W\xrightarrow{[\phi]\mapsto[\Theta\phi]}
 Q_W\longrightarrow Q\longrightarrow0,
 \qquad Q_W=\mathscr B/\Theta W.                \tag{BK.24}
\]
Injectivity of $\Theta$ proves the first map injective, and the
quotient definitions prove its image is the last map's kernel.
All maps commute with $D$ when $W$ is invariant.  Tensoring this
sequence with the flat $O$ retains the entire kernel
$O\otimes_A(V/W)$, which can have additional finite spectral
torsion before transport to the full label.  No such proper-label
torsion is discarded by the scope of the structure theorem.
The inclusion of labels $W_*\hookrightarrow V$ carries
$\mathscr B/\Theta W_*\to\mathscr B/\Theta V$ with the original
kernel $\Theta V/\Theta W_*$.  Balancing imposes the explicit
linear relations $sa\otimes v-a\otimes Dv$ within each such
label.  A linear map $L$ at a fixed label $\ell$ lifts as
$(\ell,x)\mapsto(\ell,Lx)$ and sends external $\tau$ to
external $\tau$.  Addition joins labels and applies their given
linear transports; the displayed formulas commute with those
transports.  Therefore (BK.22) is sent to the supported zero
$(\ell,0)$ by Mellin evaluation, while its nonzero kernel class
remains at that label.  Equations (BK.18)--(BK.24), including every
actual nilpotent jet and the proper-label transition kernel,
describe precisely which full-source kernel assertions survive
the complete original quotient diagram.
